\SourceLeaf{10}
\[
\mathrm{d}\Sigma=-\mathrm{d}s'.\cos QX'=+\mathrm{d}s''.\cos QX''=-\mathrm{d}s'''.\cos QX'''\text{ etc.}
\]
und daher, da die Anzahl der Teile gerade ist,
\[
\mathrm{d}s'.\cos QX'+\mathrm{d}s''.\cos QX''+\mathrm{d}s'''.\cos QX'''+\text{ etc.}=0
\]
\par
Behandeln wir alle übrigen Elemente $\mathrm{d}\Sigma$ auf dieselbe Weise und summieren, so erhalten wir
\par\medskip\Needspace{4\baselineskip}
\begin{center}ERSTER LEHRSATZ.\end{center}
\par
\textit{Das über die ganze Oberfläche des Körpers erstreckte Integral $\int\mathrm{d}s\cos QX$ ist $=0$.}
\par
Allgemeiner ergibt sich auf dieselbe Weise, dass das Integral
\[
\int(T\cos QX+U\cos QY+V\cos QZ)\mathrm{d}s
\]
verschwindet, wenn $T,U,V$ rationale Funktionen allein von $y,z$, allein von $x,z$ beziehungsweise allein von $x,y$ bezeichnen.

\Article{4}
\par
Da die Volumina der Zylinderabschnitte von unserer Ebene bis zu den Punkten $P',P'',P'''$ und so weiter der Reihe nach $=\mathrm{d}\Sigma.(x'-\alpha)$, $\mathrm{d}\Sigma.(x''-\alpha)$, $\mathrm{d}\Sigma.(x'''-\alpha)$ usw. sind, ergibt sich für den innerhalb des Zylinders liegenden Teil des Körpervolumens
\[
\begin{aligned}
&=-x'\mathrm{d}\Sigma+x''\mathrm{d}\Sigma-x'''\mathrm{d}\Sigma+\text{ etc.}\\
&=\mathrm{d}s'.x'\cos QX'+\mathrm{d}s''.x''\cos QX''+\mathrm{d}s'''.x'''\cos QX'''+\text{ etc.}
\end{aligned}
\]
und durch Summation über alle $\mathrm{d}\Sigma$ erhalten wir
\par\medskip\Needspace{4\baselineskip}
\begin{center}ZWEITER LEHRSATZ.\end{center}
\par
\textit{Das gesamte Volumen des Körpers wird durch das über die ganze Oberfläche erstreckte Integral $\int\mathrm{d}s.x\cos QX$ ausgedrückt.}
\par
Offenbar lässt sich dasselbe Volumen auch durch $\int\mathrm{d}s.y\cos QY$ oder durch $\int\mathrm{d}s.z\cos QZ$ ausdrücken.

\Article{5}
\par
Man denke sich zunächst den ganzen Zylinder mit Materie von gleichmäßiger Dichte gefüllt. Sehen wir nun, wie stark jedes seiner Elemente den Punkt $M$ \SourceLeaf{11}anzieht. Man zerlege den Zylinder durch Ebenen, die einander unendlich nahe und zur Grundfläche parallel sind, in Elementarzylinder. Ein solcher Elementarzylinder am Punkt mit den Koordinaten $\xi,\eta,\zeta$ lässt sich durch $\mathrm{d}\Sigma.\mathrm{d}\xi$ ausdrücken. Sein Abstand vom Punkt $M$ ist
\[
=\BareRoot{((a-\xi)^2+(b-\eta)^2+(c-\zeta)^2)}=\rho
\]
Damit lässt sich die Anziehung, die er auf den Punkt $M$ ausübt, durch $\mathrm{d}\Sigma.\mathrm{d}\xi.f\rho$ ausdrücken, wobei die Funktion $f\rho$ das Anziehungsgesetz bezeichnet. Da im ganzen Zylinder nur $\xi$ als veränderlich anzusehen ist, gilt $\rho\mathrm{d}\rho=-(a-\xi)\mathrm{d}\xi$. Die Anziehung des Elements ist daher $=-\frac{\rho f\rho.\mathrm{d}\rho.\mathrm{d}\Sigma}{(a-\xi)}$. Zerlegt man diese Anziehung in drei Komponenten, die zu den Koordinatenachsen $x,y,z$ parallel und ihnen entgegengerichtet sind, so ist die erste $=-f\rho.\mathrm{d}\rho.\mathrm{d}\Sigma$. Bezeichnet man nun das Integral $\int f\rho.\mathrm{d}\rho$ mit $F\rho$, so ist die Anziehung, die der von der Grundfläche $\mathrm{d}\Sigma$ bis zum Punkt mit erster Koordinate $=\xi$ reichende Zylinder auf den Punkt $M$ längs der $x$-Koordinatenachse ausübt, $=-(F\rho-\mathrm{Const.})\mathrm{d}\Sigma=-(F\rho-FR)\mathrm{d}\Sigma$, wenn $R$ den Abstand der Grundfläche $\mathrm{d}\Sigma$ vom Punkt $M$ bezeichnen soll. Folglich ist dieselbe Komponente der Anziehung aller innerhalb des Zylinders liegenden Teile des Körpers
\[
\begin{aligned}
&=(Fr'-Fr''+Fr'''-\text{ etc.})\mathrm{d}\Sigma\\
&=-Fr'.\mathrm{d}s'.\cos QX'-Fr''.\mathrm{d}s''.\cos QX''-Fr'''.\mathrm{d}s'''.\cos QX'''-\text{ etc.}
\end{aligned}
\]
\par
Wenden wir diese Überlegung auf alle Elemente $\mathrm{d}\Sigma$ an, so erhalten wir
\par\medskip\Needspace{4\baselineskip}
\begin{center}DRITTER LEHRSATZ.\end{center}
\par
{\rightskip=0pt plus 3em \textit{Die Anziehung, die der Körper auf den Punkt $M$ in der zur $x$-Koordinatenachse parallelen und ihr entgegengesetzten Richtung ausübt, wird durch das über die ganze Oberfläche erstreckte Integral $-\int Fr.\mathrm{d}s.\cos QX$ ausgedrückt.}\par}
\par
{\rightskip=0pt plus 3em Ganz entsprechend wird die Anziehung längs der beiden anderen Hauptrichtungen offenbar durch die Integrale $-\int Fr.\mathrm{d}s.\cos QY$, $-\int Fr.\mathrm{d}s.\cos QZ$ ausgedrückt.\par}

\Article{6}
\par
Nun gehen wir die Sache auf einem anderen Weg an. Man denke sich eine Kugeloberfläche mit Radius $=1$ um den Mittelpunkt $M$ beschrieben und in unendlich kleine Elemente zerlegt. Es sei $\Pi$ ein Punkt dieser Oberfläche, der zu dem kleinen Flächenstück $\mathrm{d}\Sigma$ auf ihr gehört. Man ziehe den Radius $M\Pi$ und verlängere ihn nötigenfalls unbegrenzt über die Kugeloberfläche hinaus. Es seien $P',P'',P''$ usw. die Punkte, in denen dieser Radius nacheinander die Oberfläche unseres Körpers schneidet. Dabei wird der Punkt $M$ selbst ausgeschlossen, falls er auf der Oberfläche liegt. Die Zahl dieser
